% Continuous probability: mathematical section prepared for inclusion in the main report.
% Requires amsmath, amssymb, amsthm. All analytic statements below are presented
% with explicit approximation contracts, not as already compiled Lean theorems.

\section{Continuous distributions with explicit approximation contracts}
\label{sec:continuous-laws}

A law and an integrable random variable are different objects.  A probability
law may have finite total mass and no first moment.  Accordingly, the density
construction below first produces probabilities of bounded regular events and
expectations of bounded test functions.  An expectation of an unbounded
observable is admitted only after a separate absolute-tail certificate is
provided.  This distinction includes the Cauchy law without incorrectly
assigning it an ordinary expectation.

\subsection{From a nonnegative kernel to a probability law}

A \emph{kernel presentation} consists of the following data.  The kernel
$w\colon\mathbb R\to[0,\infty)$ is given locally by primitive recursive rational
approximations, a modulus of continuity on each rational compact interval, and
a rational bound there.  Finitely many explicitly listed discontinuities are
also allowed, by splitting the interval at them.  There is a primitive
recursive radius $R(k)\geq 1$ such that, for every rational $B\geq R(k)$,
\begin{equation}
 \int_{[-B,B]\setminus[-R(k),R(k)]}w(x)\,dx\leq 2^{-k}.
 \tag{K1}\label{eq:kernel-tail}
\end{equation}
All integrals in this hypothesis are finite-interval integrals, computed by
rational quadrature with the supplied modulus.  Finally there are explicit
rationals $z_0>0$ and $R_0>0$ with
\begin{equation}
 z_0\leq\int_{-R_0}^{R_0}w(x)\,dx.
 \tag{K2}\label{eq:kernel-positive}
\end{equation}
The inequalities are finite rational approximation contracts, including any
rounding allowance; neither an unbounded search for a cutoff nor a decision
about equality of real numbers is part of the data.

\paragraph{Normalization theorem.}
These data determine a fast Cauchy name for
$Z=\int_{\mathbb R}w$, with $Z\geq z_0$, and a probability law with density
$q=w/Z$.  If $g$ is bounded by a known rational $C$ and is locally effectively
integrable against $w$, its expectation has a primitive recursive Cauchy name.
For an unbounded $g$ the same conclusion follows from a supplied primitive
recursive cutoff $R_g(k)$ for the tails of $|g|w$.

\begin{proof}
Choose a rational quadrature $a_k$ for the integral over
$[-R(k+2),R(k+2)]$ with error at most $2^{-k-2}$.  Replace $R$ by its finite
maximum up to its argument if necessary, so that it is nondecreasing.
For $\ell\geq k$, compare the two finite integrals on the larger interval.
Their difference is at most $2^{-k-2}$ by \eqref{eq:kernel-tail}.  Thus
$|a_\ell-a_k|\leq 3\cdot2^{-k-2}$.  Reindexing gives the stipulated fast
Cauchy form.  Positivity follows by retaining the fixed interval in
\eqref{eq:kernel-positive}.  Division is effective because its denominator
is separated from zero: for positive $u,v\geq z_0/2$,
\[
 |u^{-1}-v^{-1}|\leq 4z_0^{-2}|u-v|.
\]
The mass outside $[-R(k),R(k)]$ is at most $2^{-k}/z_0$.
For bounded $g$, the discarded expectation is at most $C2^{-k}/z_0$;
for unbounded $g$ use the specified $|g|w$ tail.  All remaining quantities
are finite sums.  Normalization, additivity, and positivity follow by
comparing those sums at a common requested precision.  No assertion that
every measurable function is integrable has been made.
\end{proof}

A convenient additional certificate is a local bound $q\leq K_R$ on
$[-R,R]$.  It controls threshold boundaries:
\[
 \mathbb P\{t-\eta<X\leq t+\eta\}\leq 2K_R\eta
 \qquad (|t|+\eta\leq R).
\]
Therefore bounded intervals have effective event names, and the distribution
function $F(t)=\mathbb P\{X\leq t\}$ is a represented continuous function.
If $q\leq K$ globally, $|F(s)-F(t)|\leq K|s-t|$.  These are the certificates
needed to turn continuous-test statements into assertions about exact
threshold events.

\subsection{An elementary primitive recursive exponential}

For completeness, the exponential evaluations required below can be reduced
to finite rational arithmetic.  On $0\leq t\leq T$, with $T\geq1$ an integer,
use $E_N(t)=\sum_{j=0}^{N}t^j/j!$.  After index $2T$, successive terms have
ratio at most $1/2$.  Also
$T^{2T}\leq 2^{2T(T+1)}$.  Consequently the choice
\[
 N=2T+k+2T(T+1)+2
\]
makes the remaining series tail at most $2^{-k-2}$.  The case $T=0$ is exact.
The reciprocal $1/E_N(t)$ approximates $e^{-t}$ with no greater absolute
error, because $E_N(t)\geq1$.  Multiplication, addition, and the derivative
identities for the resulting names follow from finite series identities and
these tail estimates.  The local bounds used in quadrature can thus be
proved inside the same finite arithmetic development.

The familiar estimates $e^t\geq1+t$ for $t\geq0$, $2\leq e\leq3$, and
$e^{1/2}\leq2$ have direct positive-series proofs.  They give explicit
rational estimates below.  The functions need not be introduced as an
uninterpreted analytic oracle.

\subsection{Uniform and exponential laws}

For rational $a<b$, the uniform density is
\[
 q(x)=\frac{\mathbf1_{[a,b]}(x)}{b-a}.
\]
Splitting at $a,b$ makes quadrature elementary.  Direct integration yields
\[
 \mathbb E X=\frac{a+b}{2},\qquad
 \operatorname{Var}(X)=\frac{(b-a)^2}{12},\qquad
 \mathbb E X^m=\frac{b^{m+1}-a^{m+1}}{(m+1)(b-a)}.
\]
These formulas are finite polynomial identities after denominators are
cleared.  Endpoints are null, so the choice of endpoint convention does not
affect the law.

Let $\lambda>0$ be rational, or a represented parameter with a supplied
rational lower bound $\lambda_0>0$.  The exponential density is
\[
 q_\lambda(x)=\lambda e^{-\lambda x}\mathbf1_{[0,\infty)}(x).
\]
The finite-interval fundamental theorem gives
$\int_0^R q_\lambda=1-e^{-\lambda R}$.  Since $e^{-s}\leq2^{-s}$ for integer
$s\geq0$, a cutoff $R\geq k/\lambda_0$ gives mass tail at most $2^{-k}$.
The derivative on the positive half-line is locally bounded using any
supplied upper bound on $\lambda$, and the jump at zero is a listed
breakpoint.

Integration by parts on $[0,R]$, followed by the explicit boundary and tail
bounds, gives
\begin{equation}
 \mathbb E X^m=\frac{m!}{\lambda^m},\qquad
 \mathbb E X=\lambda^{-1},\qquad
 \operatorname{Var}(X)=\lambda^{-2}.
 \label{eq:exponential-moments}
\end{equation}
Here is a uniform moment-tail certificate, so that this passage is justified
for every fixed natural $m$:
\begin{equation}
 \int_R^\infty x^m q_\lambda(x)\,dx
 \leq 2^{m+1}m!\lambda_0^{-m}e^{-\lambda_0R/2}.
 \label{eq:exponential-tail}
\end{equation}
Indeed $t^m\leq2^m m!e^{t/2}$ follows by retaining the $m$th positive term of
the exponential series.  Substitute $t=\lambda x$ and integrate the remaining
exponential.  For example, if $C$ is any integer at least
$\max(1,2^{m+1}m!\lambda_0^{-m})$, then
$R=2(k+C+1)/\lambda_0$ makes the right side at most $2^{-k}$.
This cutoff is deliberately conservative and is primitive recursive in all
displayed integer and rational inputs.

\subsection{The normal law without a hidden normalization theorem}

Put $w(x)=e^{-x^2/2}$ and define $Z$ by the normalization theorem.  On
$[-1,1]$, $w\geq1/2$, hence $Z\geq1$.  For $R\geq1$,
\begin{equation}
 \int_{|x|>R}w(x)\,dx
 \leq \frac{2}{R}e^{-R^2/2}.
 \label{eq:gaussian-tail}
\end{equation}
To prove this, first integrate only over $[R,B]$.  Since $x/R\geq1$ there,
\[
 \int_R^B e^{-x^2/2}\,dx
 \leq R^{-1}\int_R^B x e^{-x^2/2}\,dx
 \leq R^{-1}e^{-R^2/2}.
\]
Reflect and use the same finite inequality at every $B$.  In particular
$Z\leq4$.  A simple radius $R=2(k+2)$ makes the mass tail at most $2^{-k}$;
the sharper bound \eqref{eq:gaussian-tail} is usually much better.

Define the standard normal density by $\varphi=w/Z$.  It is bounded by $1$
and is globally $1$-Lipschitz.  In fact
$|w'(x)|=|x|e^{-x^2/2}\leq1$, since
$e^{x^2/2}\geq1+x^2/2\geq|x|$.
On every $[-R,R]$ it also has an explicit positive lower bound,
\begin{equation}
 \varphi(x)\geq\frac{1}{4\cdot3^{\lceil R^2/2\rceil}}.
 \label{eq:gaussian-slope}
\end{equation}
This is useful for bounded inverse-distribution calculations; it supplies
the separation that a blind bisection search would otherwise lack.

For all integers $m\geq0$ and $R\geq1$,
\begin{equation}
 \int_{|x|>R}|x|^m\varphi(x)\,dx
 \leq C_m e^{-R/4},\qquad C_m=8\cdot4^m m!.
 \label{eq:normal-moment-tail}
\end{equation}
Indeed $|x|^m\leq|x|^{2m}\leq4^m m!e^{x^2/4}$ on $|x|\geq1$;
then $e^{-x^2/4}\leq e^{-|x|/4}$ and $Z\geq1$.  Thus
$R_m(k)=4(k+C_m)$ is one primitive recursive moment cutoff, since
$C_m\leq2^{C_m}$.

Symmetry proves all odd moments zero.  Integration by parts on $[-R,R]$
and \eqref{eq:normal-moment-tail} prove
\begin{equation}
 \mathbb E G^{2m}=(2m-1)!!,\quad
 \mathbb E G=0,\quad \mathbb E G^2=1,\quad \mathbb E G^4=3.
 \label{eq:normal-moments}
\end{equation}
The convention is $(-1)!!=1$.  Explicitly, $w'=-xw$ gives the recurrence
$\int x^{j+2}w=(j+1)\int x^jw$ after the two boundary terms vanish with
the stated rate.  In particular,
\begin{equation}
 \mathbb E|G|^3\leq\tfrac12(\mathbb E G^2+\mathbb E G^4)=2,
 \label{eq:normal-third}
\end{equation}
because $2|x|^3\leq x^2+x^4$.

The usual notation $Z=\sqrt{2\pi}$ can also be justified within this
framework.  Define $\pi=4\int_0^1(1+t^2)^{-1}\,dt$.  Reciprocal substitution
shows $\int_0^\infty(1+t^2)^{-1}\,dt=\pi/2$.  In the positive quadrant,
substitute $y=tx$ in the product integral for $Z^2/4$ and then integrate in
$x$:
\[
 \frac{Z^2}{4}
 =\int_0^\infty\int_0^\infty
      x e^{-(1+t^2)x^2/2}\,dx\,dt
 =\int_0^\infty\frac{dt}{1+t^2}=\frac\pi2.
\]
This is not an invocation of an unrestricted improper change-of-variables
theorem.  Perform the substitution first on $a\leq x\leq R$, $0\leq t\leq T$.
In the original quadrant, the omitted strip $x<a$ has integral at most $2a$,
the omitted tail $x>R$ is bounded by twice the one-sided Gaussian tail, and
on $x\geq a$ the omitted region $y>Tx$ is contained in $y>Ta$.  Choose
$a,R,Ta$ with the Gaussian bounds.  In the transformed integral, the omitted
$t>T$ tail is at most $1/T$; on $t\leq T$ the omissions $x<a$ and $x>R$
are bounded by $Ta^2/2$ and $T e^{-R^2/2}$ respectively.  One simultaneous primitive recursive choice is
\[
 s=k+4,\qquad a=2^{-2s},\qquad T=2^{3s},\qquad R=2(s+1).
\]
Then $Ta=2^s$, $Ta^2/2=2^{-s-1}$, $1/T=2^{-3s}$, and
$T e^{-R^2/2}\leq2^{3s-2(s+1)^2}\leq2^{-s-2}$.
Each of the other three error terms is at most $2^{-s}$.  The sum of all
six is therefore less than $2^{-k}$.  This supplies a single explicit
finite-stage comparison rather than an unrestricted passage to an improper
integral.  Only finite Fubini and one-dimensional substitutions are used.  The normalization-free definition
$\varphi=w/Z$ and every result above are independent of this identification.

For represented $\mu$ and $\sigma\geq\sigma_0>0$, the affine image
$X=\mu+\sigma G$ has density
\[
 \varphi_{\mu,\sigma}(x)=\frac1{\sigma Z}
       \exp\!\left(-\frac{(x-\mu)^2}{2\sigma^2}\right),
\]
mean $\mu$ and variance $\sigma^2$.  The lower bound on $\sigma$ makes
division and the parameter-dependent error estimates effective.

\subsection{Cauchy: a continuous law with no expectation}

For $w(x)=(1+x^2)^{-1}$ the normalization constant again satisfies
$1\leq Z\leq4$, and reciprocal substitution gives $Z=\pi$ as defined above.
For $R\geq1$,
\[
 \int_{|x|>R}w(x)\,dx\leq 2/R.
\]
Thus $R(k)=2^{k+1}$ is a mass cutoff for the normalized density
$q(x)=1/[\pi(1+x^2)]$.  This density is bounded by $1$, is $1$-Lipschitz,
and on $[-R,R]$ is at least $1/[4(1+R^2)]$.

It has no absolute first moment.  On
$[2^j,2^{j+1}]$, $x/(1+x^2)\geq2^{-j-2}$, so the unnormalized first-moment
integral over this interval is at least $1/4$.  Reflect, sum $M$ such
intervals, and divide by $Z\leq4$ to obtain
\[
 \int_{-2^M}^{2^M}|x|q(x)\,dx\geq M/8.
\]
This finite lower bound excludes any finite absolute-moment name.  The
symmetric integrals of $xq(x)$ are zero, but that principal-value statement
is not an expectation.  A realization of the Cauchy law belongs to the
framework's convergence-in-probability class, not its $L^1$ class.

